\cleardoublepage

\section{绪论}

\subsection{问题背景}
单元测试定义->单元测试的重要性->单元测试生成的几类方法（手动编写、基于传统方法的编写，每种方法的具体方法，存在的不足）

单元测试是软件工程领域中面向软件最小可测试逻辑单元的自动化测试手段，旨在校验目标程序单元的运行行为与逻辑规约是否满足设计预期。
该测试方法通过预置测试输入、调用被测单元，并对程序输出结果及运行状态进行断言校验，可在软件开发前期及早识别功能逻辑缺陷、版本回归异常与接口适配冲突等潜在问题，
有效规避各类缺陷向后置阶段蔓延，进而大幅降低软件缺陷的修复成本，同时提升软件整体质量与代码可维护性。单元测试已成为保证软件可靠性和支持快速迭代开发的关键环节。

单元测试领域的发展经历了从手动编写到自动化生成的演进过程。最初，单元测试用例主要由开发人员或测试工程师根据需求文档和代码实现手工设计编写。
这种方式虽然能够生成可读性高、语义明确且便于维护的测试用例，但其高昂的人工成本、覆盖受限以及难以扩展到大规模代码库的缺点逐渐显现。
近年来，随着软件系统复杂度的提升和持续集成/持续部署（CI/CD）实践的普及，自动化单元测试生成方法成为研究与工业界的热点。学术界与工业界提出了多类自动化单元测试生成方法，
主要可归纳为以下几类：

（1）基于搜索/符号执行的方法：早期工作首先提出了以随机测试为核心的策略，代表性工具 Randoop 通过随机生成对象操作序列并在运行时检测异常来快速暴露运行期错误，
其实现轻量、易于集成但无法自动合成语义断言；
随后基于符号执行的 KLEE 将程序路径符号化并借助约束求解器精确生成能够触发特定路径的输入，从而在边界条件与深层路径缺陷发现上较随机方法具备更强的理论保证，
但其易遭遇路径爆炸与求解器性能瓶颈并对环境建模有较高要求；
在此基础上，后来的 EvoSuite 将搜索驱动的测试生成（以遗传算法为代表）与运行时信息相结合，提出以覆盖率或其它测试目标为适应度的演化策略，
自动生成可执行的 JUnit 测试和断言，并通过插桩、模拟（mocking）等工程化手段改善对复杂依赖的处理能力，从而在工程应用中显著提升覆盖率与断言质量。
总体来看，研究者们尝试通过更精细的路径分析、基于目标的搜索优化以及工程化的执行环境处理，
弥补早期随机生成或纯符号执行方法在断言生成、可维护性与依赖管理上的不足，但仍面临路径爆炸、断言语义准确性与跨语言/跨框架适配等挑战。

（2）基于学习与模型的方法：
于是在符号执行与搜索方法的基础上，研究者们开始引入机器学习与深度学习技术，试图通过数据驱动的方式从代码上下文中学习到更丰富的语义信息，从而生成更合理的测试输入和断言。
早期工作主要采用有监督的特征工程与浅层模型，利用手工构造的程序表示或 API 使用模式来预测测试输入或断言模板；
随后序列到序列（seq2seq）与基于编码器—解码器的神经生成模型被引入代码生成领域，使研究者能够直接从源代码或上下文学习到更丰富的语法与语义映射，
从而自动生成初步可读的测试片段。
进入预训练模型时代后，面向代码的表征模型（例如 CodeBERT、CodeT5 等）通过大规模的语料预训练并在下游任务上微调，显著提升了对代码语义的理解能力与生成质量，
能够在有限标注下生成更合理的断言与测试结构；
紧接着出现的基于大规模自回归模型的工程化产物（如 Codex、ChatGPT 等）进一步将提示工程与少样本学习结合，
支持通过自然语言提示快速生成高层次且风格多样的测试用例，增强了上下文感知能力和断言生成的语义性，
但也引入了更高的计算开销、对提示质量的敏感性以及可执行性与稳定性的问题。
为缓解这些问题，近年的研究倾向于结合检索增强生成、重排序/评分机制以及小模型与大模型的混合策略，以在保证生成语义性的同时改善可执行性和工程可用性。

尽管上述各类方法在不同维度取得了进展，但在工程化落地时仍面临一组相互关联的挑战：
一方面，为获得更高的覆盖率与更精确的断言，研究常常因为大规模模型推理而依赖计算开销大的技术，这使得在大型代码库或 CI 流水线中高频运行变得昂贵且不易扩展；
另一方面，自动生成的测试在语义正确性与可执行性之间常存在矛盾——许多方法能产生语义上合理的断言或触发边界路径的输入，
但在实际工程环境中仍需人工修补以满足可执行性与依赖管理的要求；
此外，由于随机化策略与模型概率性输出的引入，生成结果在鲁棒性与复现性上存在不确定性，增加了回归测试的维护成本；资源调度方面，
现有工作较少考虑如何在不同能力与成本的模型之间进行协同和分配，导致在可接受质量范围内仍存在资源浪费；
最后，当生成用例出现问题时，多数方法依赖整体重生成或人工修复，缺乏基于证据的局部修补与逐步改进机制。

基于上述观察，本文关注如何在不牺牲测试质量（例如断言准确性、可执行性与覆盖率）的前提下，提升生成与修复的效率并优化资源利用；
为此，我们提出一种以能力分层与证据驱动的判断器（dispatcher）为核心的工程化方案，
旨在通过小模型优先、失败证据触发升级以及闭环验证与局部修复等策略，实现高效且可控的测试生成与修复流程。
下一节“研究动机”将从基于大模型的单元测试生成出发，详细阐述方法的设计动因与整体流程。


\subsection{研究动机}
近年来，随着预训练模型在代码理解与生成任务上的成功，基于学习的方法被广泛尝试用于自动化单元测试的生成与修复，这一趋势带来了明显的工程化机遇与新的挑战。
从动机层面可以从两个角度来理解：一是质量提升的需求——基于模型的方法能够利用大规模代码与测试语料捕获复杂的语义模式，生成具备更好断言语义和更高可读性的测试片段，
从而补足传统搜索或符号执行法在断言合成与高阶语义理解上的不足；
二是工程效率的诉求——在持续集成与快速迭代的开发流程中，快速、可复用的测试生成流程能够显著降低人工负担并加速缺陷定位与回归验证。

基于模型的单元测试生成通常包含若干关键步骤：
首先是数据准备与上下文抽取，研究者需要从代码库中收集代码—测试对、提取方法签名、调用序列与依赖信息，并构造合适的输入上下文或提示模板；
其次是模型选择与训练策略，方法从早期的有监督浅层模型发展到 seq2seq 神经生成，
再到当前的预训练编码器/解码器与自回归大模型（通过微调或提示工程适配测试生成任务）；
第三是推理与生成阶段，包括温度、采样策略、束搜索与后处理（如语法校验、类型填充、断言重写）等；
最后是评估与闭环优化，常用指标涵盖执行成功率、断言准确性、代码覆盖率、发现真实缺陷的能力与生成稳定性（如非确定性或 flaky 情况的频率），
并经常结合自动化执行管道进行回归验证。

代表性工作在各环节做出不同取舍：有的侧重于通过微调提升断言生成精度，有的通过检索增强或对生成结果进行重排序以提高可执行性，还有的尝试将小模型用于快速筛选并以大模型对可疑样例进行精修以平衡成本与质量。尽管这些方法在实验室环境下取得了可观效果，但在工程化部署时仍面临成本（尤其是大模型推理成本）、可复现性、对提示及上下文依赖的敏感性，以及生成与实际工程依赖（例如第三方库、资源初始化）的对接难题。因此，本研究的动机在于设计一种面向工程化的协同生成与修复机制，通过模型能力分层、证据驱动的升级策略与闭环验证，使得在常见的 CI 场景与大规模代码库中既能利用模型的语义生成优势，又能控制资源消耗与保证生成用例的可执行性与稳定性。下一章将基于以上动机介绍本文提出的判断器（dispatcher）设计与实现要点。


\subsection{研究内容}
提出现有这些方法虽然取得了一定的进展，但在实际应用中仍然存在瓶颈和挑战：
1.挑战一是大小模型在生成单元测试代码时，性能不同，现有方法大多统一使用单一模型生成，缺乏大小模型协同生成，导致资源浪费和性能有限；
2.挑战二是在对生成的单元测试代码进行修复时，也大多采用单一模型进行修复，缺乏大小模型协同修复，导致修复过程资源浪费，修复性能提升有限。
->针对这些问题，本文提出了一个什么方法，具体是怎么做的。

\subsection{论文组织结构}
本文的组织结构如下：
第一章为绪论，主要介绍课题的研究背景、研究动机以及本文的主要研究内容，明确本文所要解决的问题与总体思路；
第二章是相关工作，回顾单元测试生成领域的主要方法与发展历程，并分析现有方法在工程化落地时面临的挑战与不足；
第三章为算法设计与实现，围绕判断器（dispatcher）机制展开，详细阐述本文所采用的分层生成策略、判定流程、反馈聚合方式以及与小模型、大模型协同工作的核心机制；
第四章为实验与评估，介绍实验数据集、对比方法、评估指标与实验设置，并对本文方法在单元测试生成任务中的有效性、资源开销与工程可用性进行系统分析。
通过以上章节安排，本文由问题提出逐步过渡到方法设计，再进一步推进到实验验证，形成完整的研究闭环。


\subsection{术语解释}
在本文中，我们将使用以下术语来描述相关概念：



